\documentclass[10pt,leqno]{amsart}
\usepackage{graphicx}
\baselineskip=16pt

\usepackage{indentfirst,csquotes}

\topmargin= .5cm
\textheight= 20cm
\textwidth= 32cc
\baselineskip=16pt

\evensidemargin= .9cm
\oddsidemargin= .9cm

\usepackage{amssymb,amsthm,amsmath}
\usepackage{xcolor,paralist,hyperref,titlesec,fancyhdr,etoolbox}
\newtheorem{theorem}{Theorem}[]
\newtheorem{definition}[theorem]{Definition}
\newtheorem{example}[theorem]{Example}
\newtheorem{lemma}[theorem]{Lemma}
\newtheorem{proposition}[theorem]{Proposition}
\newtheorem{corollary}[theorem]{Corollary}
\newtheorem{conjecture}[theorem]{Conjecture}

\titleformat{\section}[display]{\normalfont\huge\bfseries\centering}{\centering\chaptertitlename\thechapter}{10pt}{\Large}
\titlespacing*{\section}{0pt}{0ex}{0ex}

\hypersetup{ colorlinks=true, linkcolor=black, filecolor=black, urlcolor=black }
\def\proof{\noindent {\it Proof. $\, $}}
\def\endproof{\hfill $\Box$ \vskip 5 pt }

\begin{document}

\title{Deterministic Two-Zone Modeling of Transient LIBS Plasma: Correcting
the Local Thermodynamic Equilibrium Illusion via Collisional-Radiative
Matrices and Statistical Inversion}

\author[Walidain, Idris, Saddami]{Birrul Walidain, Nasrullah Idris, Khairun Saddami}
\date{\today}
\address{Department of Physics, Faculty of Mathematics and Natural Sciences,
Universitas Syiah Kuala, Banda Aceh, Indonesia}
\email{birrul.walidain@mhs.usk.ac.id}
\maketitle

\let\thefootnote\relax
\footnotetext{MSC2020: Primary 82D10 (Plasma and Plasmonics), Secondary 65L04 (Numerical Methods for Stiff ODE), 62J02 (General Nonlinear Regression).}

\begin{abstract}
The extraction of fundamental thermodynamic parameters, such as electron
temperature ($T$) and density ($n_e$), in Laser-Induced Breakdown Spectroscopy
(LIBS) frequently relies on the assumption of Local Thermodynamic Equilibrium
(LTE) via the Saha equation. However, in rapidly expanding transient plasmas
operating on nanosecond scales, the LTE criterion collapses, and intense spatial
gradients induce severe optical self-absorption. This paper presents a
deterministic computational architecture that replaces the homogeneous plasma
assumption with a Two-Zone Model (Core and Shell). We strictly reject the Saha
equilibrium assumption, instead solving the highly stiff Ordinary Differential
Equations (ODEs) of a Collisional-Radiative (CR) model using an implicit solver
with an analytical Jacobian matrix to preserve absolute thermodynamic
conservation. Furthermore, rather than employing computationally expensive
stochastic blind-search algorithms or black-box deep learning models, we utilize
a statistical machine learning regression trained on deterministically
synthesized spectra. This inverse mathematical solver effectively deconvolutes
the self-absorbed macroscopic spectra to reveal the true microscopic kinetic
parameters of the core, thereby correcting the McWhirter criterion and
eliminating the LTE illusion.
\end{abstract}

\bigskip

% ============================================================================
\section{Introduction}
% ============================================================================

Laser-Induced Breakdown Spectroscopy (LIBS) has been widely established as a
robust analytical technique for in-situ material characterization. Despite its
empirical success, the theoretical extraction of plasma parameters relies heavily
on the assumption that the plasma adheres to Local Thermodynamic Equilibrium
(LTE). Under LTE, atomic state populations are assumed to follow the Boltzmann
distribution, and ionization balances are governed by the Saha equation.
However, laser-induced plasmas are inherently transient and chaotic, exhibiting
extreme spatial gradients and rapid expansion velocities on nanosecond and
femtosecond timescales. In such regimes, collisional rates often fail to
equilibrate with radiative decay processes, causing a fundamental breakdown of
the LTE condition.

Furthermore, the sharp temperature gradient between the hot, dense core and the
cooler, expanding periphery induces profound optical self-absorption. Recent
literature has attempted to address this self-reversal phenomenon by applying
homogeneous plasma models or by utilizing global stochastic optimization
algorithms (such as \texttt{CRS2-LM}) to fit experimental spectra. These
approaches contain fatal thermodynamic flaws: they either ignore the spatial
non-uniformity of the transient expansion or forcibly embed the Saha equilibrium
assumption into mathematically partitioned spatial zones \cite{griem1997,
fujimoto2004}.

In rejecting black-box deep learning architectures and probabilistic
curve-fitting that obscure physical realities, this study proposes a strictly
deterministic computational framework. We implement a physically rigorous
Two-Zone Model where the atomic population kinetics are governed purely by an
exact Collisional-Radiative (CR) matrix, eliminating the Saha assumption
entirely. By mapping the deterministic forward model to a rapid statistical
machine learning regression for spectral inversion, this methodology provides a
high-fidelity diagnostic tool capable of exposing the LTE illusion and
validating the empirical limits of the McWhirter criteria \cite{mcwhirter1965}.

% ============================================================================
\section{Material and Methods}
% ============================================================================

The computational architecture is strictly divided into a deterministic forward
physical model---ensuring absolute conservation of mass and energy---and a
statistical inverse solver for parameter extraction.

\subsection{The Two-Zone Geometric Architecture}

To capture the spatial heterogeneity of the self-absorption phenomenon without
triggering the exponential computational cost of full 3D hydrodynamic
simulations, the plasma is spatially discretized into two distinct regions: the
emissive Core ($V_1$) and the absorptive Shell ($V_p$). The macroscopic intensity
observed by the spectrometer ($I_\mathrm{obs}$) is modeled via the 1D integration
of the Radiative Transfer Equation (RTE):
\begin{equation}
    I_\mathrm{obs}(\lambda) = I_\mathrm{core}(\lambda)\,
    \exp\!\bigl(-\kappa_\mathrm{shell}(\lambda)\, d_\mathrm{shell}\bigr)
    + I_\mathrm{shell}(\lambda),
    \label{eq:rte}
\end{equation}
where $I_\mathrm{core}$ is the unattenuated emission from the core,
$\kappa_\mathrm{shell}$ is the wavelength-dependent absorption coefficient of
the periphery, and $d_\mathrm{shell}$ is the optical thickness of the absorbing
layer \cite{rybicki1979}.

\subsection{Deterministic Population Kinetics (The Core)}

The fundamental divergence of this methodology from established homogeneous
models lies in the calculation of the energy state populations ($n_i$). We
unequivocally reject the Saha equation. The temporal evolution of the fractional
populations is solved via the Collisional-Radiative (CR) system of differential
equations:
\begin{equation}
    \frac{dn_i}{dt} = \sum_{j \neq i} n_j\, R_{ji} - n_i \sum_{j \neq i} R_{ij},
    \label{eq:cr}
\end{equation}
where the rate matrix $R_{ij}$ encapsulates all fundamental electron-impact
collisional rates and radiative Einstein transition probabilities ($A_{ij}$,
$B_{ij}$). Because the kinetic rates of collisional excitation and spontaneous
emission are separated by multiple orders of magnitude (ranging from $10^{-14}$
to $10^{-3}$ seconds), the resulting Jacobian matrix is exceptionally stiff.

To prevent numerical divergence and unphysical negative particle
populations---which routinely occur when using explicit methods such as standard
Runge-Kutta---the temporal integration is executed using an implicit solver
(Backward Differentiation Formula, BDF, or Radau). The computational bottleneck
of inverting the $\mathcal{O}(N^3)$ matrix at each step is mitigated by
supplying the exact analytical Jacobian:
\begin{equation}
    J_{ij} = \frac{\partial f_i}{\partial n_j},
    \label{eq:jacobian}
\end{equation}
which is analytically identical to the CR rate matrix $M$ for the linear ODE
system $\mathbf{f}(n) = M\,n$. This guarantees that Newton iterations within
the implicit solver converge quadratically, preserving strict thermodynamic
conservation at every time step \cite{hairer1996}.

\subsection{Socratic Boundaries and Synthetic Dataset Generation}

The forward model computes the absorption coefficients ($\kappa$) by summing
Voigt line profiles over all NIST-validated transitions, integrating both Stark
broadening (pressure broadening, Lorentzian) and Doppler broadening (thermal
motion, Gaussian):
\begin{equation}
    \kappa(\lambda) = \sum_{\text{transitions}} n_i\,\sigma_{ij}(\lambda)\,
    \left(1 - \frac{n_k\,g_i}{n_i\,g_k}\right),
    \label{eq:kappa}
\end{equation}
where the stimulated emission correction in parentheses ensures detailed balance
is maintained \cite{griem1997}.

By sweeping the parameter space ($T_e \in [4000, 20000]$\,K and
$n_e \in [10^{14}, 10^{18}]$\,cm$^{-3}$), we synthesize a comprehensive dataset
of self-absorbed spectra. We apply Socratic ignorance regarding fundamental
atomic inputs: transition probabilities $A_{ki}$ and statistical weights $g_i$
are strictly sourced from the NIST Atomic Spectra Database \cite{nist_asd}.
Transitions lacking validated empirical cross-sections are not synthetically
approximated, ensuring the mathematical space does not exceed the boundaries of
verified physical literature.

\subsection{Inverse Modeling via Statistical Regression}

To solve the inverse problem---extracting the pure core temperature
($T_\mathrm{core}$) and electron density ($n_{e,\mathrm{core}}$) from a
convoluted experimental spectrum---we bypass computationally exhaustive and
often divergent stochastic blind-search algorithms. Instead, intermediate
statistical machine learning algorithms (e.g., Support Vector Regression or
Ridge Regression) are trained on the deterministically generated synthetic
dataset. Here, the algorithm acts purely as a high-speed mathematical inversion
operator rather than a predictive physical model, translating the distorted line
profiles back into their foundational thermodynamic parameters.

% ============================================================================
\section{Expected Results and Discussion}
% ============================================================================

By applying the trained regression model to experimental LIBS spectra, this
computational architecture is expected to successfully deconvolute the optical
self-absorption. Consequently, the extracted parameters ($T$ and $n_e$) will
reflect the true, unattenuated state of the plasma core.

When these corrected parameters are substituted back into the classical
McWhirter criterion \cite{mcwhirter1965}:
\begin{equation}
    N_e \ge 1.6 \times 10^{12}\, T^{1/2}\, (\Delta E)^3,
    \label{eq:mcwhirter}
\end{equation}
we hypothesize a significant deviation from the values derived via traditional
homogeneous Boltzmann plot methods. This will quantitatively demonstrate that
previous assertions of LTE in transient LIBS plasmas are largely mathematical
illusions caused by the failure to isolate the absorptive shell from the
emissive core. The stability of the implicit CR solver guarantees that this
correction is grounded in absolute thermodynamic conservation rather than
statistical curve-fitting anomalies.

% ============================================================================
\begin{thebibliography}{99}

\bibitem{fujimoto2004}
Fujimoto, T. (2004).
\textit{Plasma Spectroscopy}.
Clarendon Press, Oxford.

\bibitem{griem1997}
Griem, H. R. (1997).
\textit{Principles of Plasma Spectroscopy}.
Cambridge University Press, Cambridge.

\bibitem{hairer1996}
Hairer, E., \& Wanner, G. (1996).
\textit{Solving Ordinary Differential Equations II: Stiff and
Differential-Algebraic Problems} (2nd ed.).
Springer, Berlin.

\bibitem{mcwhirter1965}
McWhirter, R. W. P. (1965).
Spectral intensities.
In R. H. Huddlestone \& S. L. Leonard (Eds.),
\textit{Plasma Diagnostic Techniques} (pp.\ 201--264).
Academic Press, New York.

\bibitem{nist_asd}
Kramida, A., Ralchenko, Yu., Reader, J., \& NIST ASD Team (2023).
\textit{NIST Atomic Spectra Database} (version 5.11).
National Institute of Standards and Technology, Gaithersburg, MD.
\url{https://physics.nist.gov/asd}

\bibitem{rybicki1979}
Rybicki, G. B., \& Lightman, A. P. (1979).
\textit{Radiative Processes in Astrophysics}.
Wiley-Interscience, New York.

\bibitem{vanregemorter1962}
van Regemorter, H. (1962).
Rate of collisional excitation in stellar atmospheres.
\textit{The Astrophysical Journal}, \textbf{136}, 906--915.

\end{thebibliography}

\end{document}